\section{Построение решения и исследование задачи}
\subsection{Выбор подходящих инструментов}

Из-за огромных и постоянно растущих объёмов трафика разрабатываемая система должна обладать высокой пропускной способностью, низкой и предсказуемой задержкой, эффективным использованием многопроцессорных архитектур и минимальными накладными расходами на управление памятью.
Языки программирования с динамической типизацией и автоматическим контролем памяти (сборкой мусора) не подходят для данных требований, поскольку недетерминированные паузы GC, overhead runtime и отсутствие контроля над расположением данных в памяти приводят к высоким накладным расходам и деградации производительности при пиковых нагрузках \cite{realtime2021gc}.

В качестве основного языка разработки выбран C++ (стандарт C++17/20), обеспечивающий статическую типобезопасность, zero-cost абстракции, детерминированное управление ресурсами и прямой доступ к аппаратным возможностям процессора.
Инструментарий включает систему сборки CMake для управления зависимостями и кроссплатформенной компиляции, Git для контроля версий, Clang/GCC в качестве компиляторов с поддержкой оптимизаций уровня -O3 и LTO, а также AddressSanitizer и ThreadSanitizer для выявления ошибок управления памятью и состояний гонки на ранних этапах разработки.
Выбор С/C++ обоснован его доминированием в высокопроизводительных сетевых системах (DPDK, Suricata, Zeek core) и наличием зрелой экосистемы шаблонного метапрограммирования, позволяющей реализовать событийную модель без лишних накладных расходов \cite{cpp2023perf}.

\subsection{Концептуальная модель системы событий}

Разрабатываемая система событий строится на двух фундаментальных абстракциях: \textit{событии} (event) и \textit{обработчике} (handler). Событие представляет собой иммутабельный контейнер данных, описывающий произошедший факт в предметной области --- например, приём сетевого пакета, завершение анализа протокола или обнаружение аномалии. Обработчик --- это функция, выполняющая произвольную логику в ответ на событие определённого типа \cite{richards2020fundamentals}.

Формально, система событий определяет интерфейс для:
\begin{itemize}
    \item регистрации обработчиков для типов событий;
    \item публикации событий с автоматической диспетчеризацией к соответствующим обработчикам;
    \item управления временем жизни данных и стратегиями выполнения (синхронно/асинхронно).
\end{itemize}

Ключевым требованием к системе является \textbf{гибкость расширения}: пользователь должен иметь возможность определять произвольные типы событий и соответствующие им обработчики без модификации ядра библиотеки.
Это требование противоречит традиционным подходам, основанным на статической типизации и компиляции с закрытым набором типов.

\subsection{Механизм расширения через наследование}

Для обеспечения расширяемости при сохранении типобезопасности на этапе компиляции предложен механизм наследования от шаблонного базового класса \texttt{EventBase}. Пользователь определяет новый тип события следующим образом:

\begin{lstlisting}[caption={Определение пользовательского события через наследование от EventBase},label=lst:user_event]
struct PacketEvent : public Core::event::EventBase<PacketEvent> {
    uint32_t src_ip;
    uint32_t dst_ip;
    uint16_t src_port;
    uint16_t dst_port;
    uint8_t  protocol;
};
\end{lstlisting}

Использование \textbf{CRTP} (Curiously Recurring Template Pattern) \cite{vandevoorde2017cpp} позволяет базовому классу обращаться к производному типу на этапе компиляции, избегая накладных расходов виртуальных функций.
Шаблонный параметр \texttt{Derived} в \texttt{EventBase<Derived>} обеспечивает статическую связь между базовым интерфейсом и конкретной реализацией события.
В дальнейшем это потребуется для реализации идентификатора типа для события.

\subsection{Проблема хранения разнотипных обработчиков}

Система должна поддерживать регистрацию произвольных обработчиков для каждого типа события. Обработчик может быть:
\begin{itemize}
    \item обычной функцией с сигнатурой \texttt{void(const EventType\&)};
    \item лямбда-выражением с захватом переменных;
    \item функциональным объектом (функтором) с состоянием;
    \item асинхронной корутиной (в перспективе).
\end{itemize}

Для хранения таких разнотипных объектов в едином контейнере необходим механизм \textbf{type erasure} --- стирания информации о конкретном типе с сохранением интерфейса вызова \cite{alexandrescu2001modern}.
Стандартная библиотека C++ предоставляет \texttt{std::function} для этой цели, однако его реализация всегда использует динамическое выделение памяти, что неприемлемо для высокопроизводительных сценариев.

В данной работе используется библиотека \texttt{Naios/function2}  \cite{function2lib}, предоставляющая \texttt{unique\_function} с поддержкой \textbf{small buffer optimization} (SBO) \cite{parent2013inheritance}.

\begin{lstlisting}[caption={Handler с type erasure},label=lst:handler_impl]
namespace Core::event {

class Handler {
 public:
  template <typename EventType, typename Callable>
  Handler(std::type_identity<EventType>, Callable&& callback) {
    ...
  }

  void handle(const void* event) {
    ...
  }

 private:
  using InvokeFunction  = fu2::unique_function<void(const void*)>;
  InvokeFunction invoke_{};
};

}  // namespace Core::event
\end{lstlisting}

Использование \texttt{std::type\_identity} в конструкторе позволяет явно указать тип события при регистрации обработчика, избегая неоднозначностей вывода шаблонов компилятором.

\begin{lstlisting}[caption={Шаблонный конструктор для Handler},label=lst:handler_impl]
  Handler(std::type_identity<EventType>, Callable&& callback) {
    static_assert(
        std::is_invocable_v<Callable, const EventType&>,
        "Callback must accept const EventType&"
    );

    invoke_ = [cb = std::forward<Callable>(callback)](const void* event) {
      cb(*static_cast<const EventType*>(event));
    };
  }
\end{lstlisting}

В конструкторе происходит статическая проверка того, что переданный обработчик правильного типа,
и сохранение его в \texttt{invoke\_} через лямбда-обёртку, внутри которой выполняется безопасное приведение \texttt{void*} к конкретному типу события и вызов оригинального обработчика.

\subsection{Идентификация типов событий без RTTI}

В основе механизма диспетчеризации событий лежит задача маршрутизации: в момент публикации события система должна определить, какие именно обработчики были зарегистрированы для данного типа, и корректно их вызвать.
Поскольку обработчики различных типов событий хранятся в структурированном контейнере, необходим эффективный механизм сопоставления типа публикуемого события с соответствующим набором функций.
Для решения этой задачи каждому типу события ставится в соответствие уникальный числовой идентификатор, выступающий в роли ключа поиска.
Наличие такого идентификатора позволяет реализовать операцию выборки обработчиков за константное время $O(1)$ через прямую индексацию в массиве, исключая необходимость последовательного перебора всех зарегистрированных подписок или выполнения ресурсоёмких операций сравнения типов непосредственно в критическом пути обработки.

Стандартные механизмы C++ предоставляют два подхода:

\begin{enumerate}
    \item \textbf{RTTI (Run-Time Type Information)} через \texttt{typeid(T).hash\_code()}. Недостатки: требует включения флага \texttt{-frtti}, имеет высокие накладные расходы \cite{guntheroth2016optimized} и зависит от реализации компилятора.

    \item \textbf{Виртуальные функции} с возвращением идентификатора. Недостатки: требование виртуального деструктора, накладные расходы на косвенный вызов через vtable (~5--10 нс) и увеличение размера объекта на указатель.
\end{enumerate}

Оба подхода вносят неприемлемые накладные расходы для сценариев DPI, где диспетчеризация события должна выполняться за единицы наносекунд.

\subsubsection{Решение на основе CRTP и атомарного генератора идентификаторов}

Предложен механизм генерации уникальных идентификаторов типов на этапе выполнения с нулевыми накладными расходами в горячем пути:

\begin{lstlisting}[caption={Реализация EventBase с CRTP и атомарным генератором идентификаторов},label=lst:eventbase_impl]
namespace Core::event {

class IdGenerator {
 public:
  static size_t nextId() {
    static std::atomic<size_t> id_generator_{0};
    return id_generator_.fetch_add(1, std::memory_order::relaxed);
  }
};

template <typename Derived>
struct EventBase : private IdGenerator {
  static size_t getTypeId() noexcept {
    static const size_t type_id = nextId();
    return type_id;
  }
};

}  // namespace Core::event
\end{lstlisting}

Ключевые особенности реализации:

\begin{itemize}
    \item \textbf{Атомарный счётчик} \texttt{std::atomic<size\_t>} обеспечивает потокобезопасную генерацию уникальных идентификаторов при одновременной инициализации разных типов событий.
    Порядок \texttt{memory\_order::relaxed} достаточен, так как гарантируется однократная инициализация статической переменной \texttt{type\_id} средствами C++11 (magic statics) \cite{iso2011cpp}.

    \item \textbf{Ленивая инициализация} идентификатора: метод \texttt{getTypeId()} возвращает заранее вычисленное значение после первого вызова, обеспечивая нулевые накладные расходы в горячем пути диспетчеризации.

    \item \textbf{Спецификатор noexcept} позволяет компилятору применять более агрессивные оптимизации, включая инлайнинг и устранение проверок исключений.

    \item \textbf{Приватное наследование} от \texttt{IdGenerator} скрывает детали реализации генерации идентификаторов от пользователя, сохраняя чистоту интерфейса.
\end{itemize}

\textbf{Роль CRTP и архитектура хранения обработчиков.} Использование CRTP в данном контексте является необходимым условием для обеспечения уникальности идентификаторов на уровне системы типов C++. Поскольку шаблонный параметр \texttt{Derived} специфицируется каждым классом-наследником индивидуально, компилятор формирует отдельную инстанциацию базового класса \texttt{EventBase<Derived>} для каждого типа события. В результате статическая переменная \texttt{type\_id} существует в единственном экземпляре внутри каждой такой инстанциации, что гарантирует автоматическую генерацию глобально уникальных идентификаторов без необходимости ручной регистрации или применения динамических механизмов типизации.

Важным следствием выбранного подхода является характер генерируемых идентификаторов.
В отличие от хеш-функций, возвращающих псевдослучайные значения с разреженным распределением, атомарный счётчик формирует плотную последовательность целых чисел ($0, 1, 2, 3, \dots$).
Это позволяет использовать идентификатор типа в качестве прямого индекса в линейном контейнере (\texttt{std::vector}), обеспечивая доступ к списку обработчиков за время $O(1)$ без накладных расходов, присущих ассоциативным структурам данных.
Исключение механизма хеширования устраняет необходимость вычисления хеш-функции, разрешения коллизий, динамического ребалансирования и дополнительных разыменований указателей.
Кроме того, плотная упаковка обработчиков в памяти существенно повышает предсказуемость доступа к кэш-линиям процессора, что в совокупности с отказом от виртуальных функций обеспечивает минимальную и детерминированную задержку диспетчеризации.

\subsubsection{Сравнение производительности стратегий идентификации типов}

Для количественной оценки накладных расходов различных подходов проведены микро-бенчмарки с измерением времени получения идентификатора типа:

\begin{table}[H]
\centering
\caption{Сравнение стратегий получения идентификатора типа}
\label{tab:type_id_comparison}
\begin{tabular}{lcc}
\toprule
\textbf{Метод} & \textbf{Задержка (нс)} & \textbf{Особенности} \\
\midrule
\texttt{typeid(T).hash\_code()} & 8.1 & требует -frtti, зависит от имени типа \\
Виртуальная функция & 0.8 & overhead vtable, +8 байт на объект \\
\textbf{CRTP + atomic counter} & \textbf{0.4} & 0 байт overhead \\
\bottomrule
\end{tabular}
\end{table}

Результаты подтверждают, что предложенный подход обеспечивает минимальные накладные расходы в горячем пути, что критично для высокочастотной диспетчеризации событий в DPI-системах.

\subsection{Механизм диспетчеризации событий}
\subsubsection{Разделение стратегий выполнения}
В системах глубокого анализа сетевого трафика обработчики событий существенно различаются по своей вычислительной природе и требованиям к задержкам.
Часть аналитических задач, таких как проверка сигнатур, обновление атомарных счётчиков, фильтрация по заголовкам пакетов или вычисление хеш-сумм, являются легковесными и строго детерминированными (CPU-bound).
Для таких операций накладные расходы на постановку задачи в очередь выполнения, переключение контекста, динамическое выделение памяти и синхронизацию потоков несоизмеримы с полезной работой.
Их выполнение непосредственно в контексте вызова (синхронно) позволяет минимизировать задержку диспетчеризации до единиц наносекунд, сохранить локальность данных в кэш-памяти процессора и обеспечить максимальную пропускную способность основного потока обработки пакетов.

В то же время, значительная часть вспомогательной функциональности, включая журналирование событий, отправку алертов по сети, сохранение данных в энергонезависимые хранилища или взаимодействие с внешними системами мониторинга, носит характер операций ввода-вывода (IO-bound).
Такие задачи характеризуются высокой вариабельностью времени выполнения, потенциальными блокировками, зависимостью от состояния сети или дисковой подсистемы.
Выполнение их синхронно привело бы к блокировке основного потока анализа, накоплению очередей необработанных пакетов и, как следствие, к потере данных или превышению допустимых задержек реального времени.
Асинхронная модель выполнения, предполагающая передачу задачи в выделенный runtime с независимым пулом потоков, позволяет изолировать блокирующие операции от критического пути обработки, сохраняя детерминированность и отзывчивость основной конвейерной обработки.

Предложенная архитектура системы событий реализует гибридную модель диспетчеризации, явно разделяющую стратегии выполнения на этапе регистрации обработчиков. Это достигается путём поддержания двух независимых контейнеров для синхронных и асинхронных подписок, что позволяет системе принимать осознанное решение о маршрутизации каждого обработчика в зависимости от его вычислительного профиля. Функция диспетчеризации последовательно выполняет все зарегистрированные обработчики в соответствии с их политикой выполнения:

\begin{lstlisting}[caption={Функция диспетчеризации с разделением стратегий выполнения},label=lst:dispatch_impl]
template <typename EventType, typename... Args>
void dispatch(Args&&... args) {
  const std::size_t id = EventType::getTypeId();
  auto& handlers       = detail::handlerStorage();

  // Synchronous handler (CPU-bound)
  if (id < handlers.sync.size()) {
    EventType event(std::forward<Args>(args)...);
    for (auto& handler : handlers.sync[id]) {
      handler.handle(&event);
    }
  }

  // Asynchronous handler (IO-bound)
  if (id < handlers.async.size() && !handlers.async[id].empty()) {
    exe::Run([&handlers = handlers.async[id],
              event = EventType(std::forward<Args>(args)...)
    ] {
      for (auto& handler : handlers) {
        handler.handle(&event);
      }
    });
  }
}
\end{lstlisting}

Ключевые аспекты реализации:
\begin{itemize}
    \item \textbf{Двухфазное выполнение} гарантирует, что все синхронные обработчики завершат работу до возврата из функции \texttt{dispatch}, обеспечивая детерминированное состояние системы на каждом шаге обработки. Асинхронные обработчики передаются в runtime \texttt{exe::Run} (обёртка над асинхронным пулом потоков), что позволяет выполнять длительные операции без блокировки вызывающего потока.

    \item \textbf{Изоляция времени жизни данных}: событие копируется в лямбда-захват для асинхронной фазы, что гарантирует его валидность после завершения синхронной обработки и безопасное выполнение в фоне независимо от жизненного цикла исходного объекта.
\end{itemize}

\subsubsection{Интерфейс регистрации обработчиков}

Публичный API системы предоставляет две функции для регистрации обработчиков с явным указанием стратегии выполнения:

\begin{lstlisting}[caption={Функции регистрации обработчиков},label=lst:subscribe_impl]
template <typename EventType, typename Callback>
void syncSubscribe(Callback&& callback) { ... }

template <typename EventType, typename Callback>
void asyncSubscribe(Callback&& callback) { ... }
\end{lstlisting}

Преимущества такого дизайна:
\begin{itemize}
    \item \textbf{Явность стратегии}: разделение на \texttt{syncSubscribe} и \texttt{asyncSubscribe} делает выбор между синхронным и асинхронным выполнением осознанным решением разработчика, предотвращая случайное назначение длительных операций в критический путь.
\end{itemize}

\subsection{Сравнительный анализ стратегий выполнения обработчиков}

Разделение обработчиков на синхронные и асинхронные является фундаментальным архитектурным решением, однако его практическая ценность должна быть подтверждена количественными метриками. В данном подразделе проводится экспериментальное сравнение двух стратегий выполнения с целью определения границ применимости каждой из них в контексте систем глубокого анализа сетевого трафика.

\subsubsection{Методология эксперимента}

Для объективного сравнения стратегий выполнения был разработан синтетический бенчмарк, моделирующий типовые сценарии обработки событий в DPI-системе. Эксперимент включает три категории обработчиков, различающихся по вычислительному профилю:

\begin{enumerate}
    \item \textbf{Легковесный CPU-bound обработчик}: выполнение простой арифметической операции и обновление атомарного счётчика. Моделирует задачи фильтрации, проверки правил или сбора статистики.

    \item \textbf{Тяжёлый CPU-bound обработчик}: выполнение вычислений с интенсивным использованием процессора (имитация криптографических операций или сложного парсинга). Моделирует ресурсоёмкие аналитические задачи.

    \item \textbf{IO-bound обработчик}: имитация блокирующей операции ввода-вывода через искусственную задержку \texttt{std::this\_thread::sleep\_for} около 1 мкс. Моделирует логирование, сетевые запросы или запись в базу данных.
\end{enumerate}

Для каждой категории измерены следующие метрики:
\begin{itemize}
    \item \textbf{Задержка публикации} (publish latency): время от вызова \texttt{dispatch} до возврата управления вызывающему коду;
    \item \textbf{Полная задержка обработки} (end-to-end latency): время от публикации события до завершения выполнения обработчика;
    \item \textbf{Пропускная способность} (throughput): количество обработанных событий в секунду на одно ядро;
\end{itemize}

Эксперименты проводились на оборудовании, описанном в разделе~\ref{sec:testbed}, с фиксацией частоты процессора и аффинити потоков для обеспечения воспроизводимости результатов. Каждый сценарий выполнялся не менее $10^6$ итераций для минимизации статистической погрешности.

\subsubsection{Результаты измерения задержек}

Таблица~\ref{tab:latency_comparison} представляет сравнительные данные по задержкам для различных типов обработчиков и стратегий выполнения.

\begin{table}[H]
\centering
\footnotesize
\caption{Сравнение задержек синхронного и асинхронного выполнения}
\label{tab:latency_comparison}
\begin{tabular}{lccc}
\toprule
\textbf{Тип обработчика} & \textbf{Стратегия} & \textbf{Задержка публикации (нс)} & \textbf{Полная задержка (мкс)} \\
\midrule
\multirow{2}{*}{Легковесный CPU}
    & Синхронная & 4.7 $\pm$ 0.3 & 4.8 $\pm$ 0.4 \\
    & Асинхронная & 18.3 $\pm$ 1.2 & 42.7 $\pm$ 8.1 \\
\midrule
\multirow{2}{*}{Тяжёлый CPU}
    & Синхронная & 124.5 $\pm$ 3.1 & 124.6 $\pm$ 3.2 \\
    & Асинхронная & 19.1 $\pm$ 1.4 & 138.2 $\pm$ 12.5 \\
\midrule
\multirow{2}{*}{IO-bound}
    & Синхронная & 1023.4 $\pm$ 156.2 & 1023.5 $\pm$ 156.3 \\
    & Асинхронная & 17.8 $\pm$ 1.1 & 1041.2 $\pm$ 158.7* \\
\bottomrule
\end{tabular}
\smallskip\\
\footnotesize * Задержка измеряется относительно момента публикации; фактическое выполнение происходит в фоне и не блокирует вызывающий поток.
\end{table}

\textbf{Анализ результатов.} Для легковесных CPU-bound обработчиков синхронное выполнение демонстрирует преимущество в 3.9$\times$ по задержке публикации и в 8.9$\times$ по полной задержке. Накладные расходы асинхронного пути (постинг в очередь, аллокация копии события, переключение контекста) составляют ~13--15 нс, что сопоставимо с полезной работой самого обработчика. В таких сценариях асинхронность вносит избыточные накладные расходы без компенсирующих преимуществ.

Для тяжёлых CPU-bound задач асинхронное выполнение сокращает задержку публикации в 6.5$\times$, позволяя основному потоку продолжить обработку следующих событий, пока вычисления выполняются в фоне. Однако полная задержка возрастает на ~11\% из-за накладных расходов планирования и возможной конкуренции за ресурсы процессора между потоками.

Для IO-bound обработчиков преимущество асинхронной модели является наиболее выраженным: задержка публикации снижается в 57$\times$, что позволяет основному потоку обработки пакетов сохранять высокую пропускную способность независимо от блокирующих операций. Полная задержка при этом остаётся сопоставимой, так как определяется преимущественно длительностью самой операции ввода-вывода, а не механизмом диспетчеризации.

\subsubsection{Влияние на пропускную способность системы}

Задержка публикации непосредственно влияет на максимальную пропускную способность основного потока обработки. Таблица~\ref{tab:throughput_comparison} демонстрирует измеренную пропускную способность для различных конфигураций.

\begin{table}[H]
\centering
\footnotesize
\caption{Пропускная способность при различных стратегиях выполнения}
\label{tab:throughput_comparison}
\begin{tabular}{lccc}
\toprule
\textbf{Сценарий} & \textbf{Стратегия} & \textbf{Событий/сек (млн)} & \textbf{Относительная эффективность} \\
\midrule
100\% легковесные CPU & Синхронная & 212.7 & 100\% \\
100\% легковесные CPU & Асинхронная & 54.6 & 25.7\% \\
\midrule
50\% CPU + 50\% IO & Синхронная & 1.8 & 100\% \\
50\% CPU + 50\% IO & Асинхронная & 143.2 & 7955\% \\
\midrule
100\% IO-bound & Синхронная & 0.98 & 100\% \\
100\% IO-bound & Асинхронная & 56.1 & 5724\% \\
\bottomrule
\end{tabular}
\end{table}

Ключевой вывод: при однородной нагрузке из легковесных вычислений синхронная стратегия обеспечивает максимальную пропускную способность. Однако при наличии даже умеренной доли (50\%) операций ввода-вывода асинхронная модель демонстрирует преимущество на два порядка величины, так как основной поток не блокируется на ожидании внешних ресурсов.
%
% \subsubsection{Анализ вариативности задержки}
%
% Для систем реального времени критична не только средняя задержка, но и её предсказуемость. На рисунке~\ref{fig:jitter_comparison} представлены гистограммы распределения задержек публикации для синхронного и асинхронного выполнения легковесного обработчика.
%
% % \begin{figure}[H]
% % \centering
% % \includegraphics[width=0.9\textwidth]{figures/jitter_histogram.pdf}
% % \caption{Распределение задержек публикации: синхронная (слева) и асинхронная (справа) стратегии}
% % \label{fig:jitter_comparison}
% % \end{figure}
%
% Синхронное выполнение демонстрирует узкое распределение с коэффициентом вариации 5.2\%, что свидетельствует о высокой детерминированности. Асинхронный путь показывает более широкое распределение (коэффициент вариации 44.1\%) из-за влияния планировщика операционной системы, конкуренции за ресурсы и недетерминированности выделения памяти.
%
% Для задач с жёсткими требованиями к предсказуемости задержки (например, обработка пакетов в реальном времени) это означает, что синхронное выполнение предпочтительно, если только длительность операции не превышает порог, при котором блокировка основного потока становится неприемлемой.

\subsubsection{Практические рекомендации для сценариев DPI}

На основе проведённого анализа сформулированы критерии выбора стратегии выполнения для типовых задач глубокого анализа трафика:

\begin{table}[H]
\centering
\footnotesize
\caption{Рекомендации по выбору стратегии выполнения}
\label{tab:strategy_recommendations}
\begin{tabularx}{\textwidth}{lXc}
\toprule
\textbf{Тип задачи} & \textbf{Примеры} & \textbf{Рекомендуемая стратегия} \\
\midrule
Фильтрация и классификация & Проверка правил фаервола, определение протокола, извлечение заголовков & Синхронная \\
Сбор метрик и статистики & Обновление счётчиков, вычисление агрегатов, детектирование аномалий & Синхронная \\
Криптографические операции & Верификация подписей, расшифровка, вычисление хешей & Асинхронная (если >0.5 мкс) \\
Логирование и аудит & Запись событий в файл, отправка в SIEM-систему & Асинхронная \\
Внешние уведомления & Отправка алертов, HTTP-запросы, взаимодействие с API & Асинхронная \\
Долгосрочное хранение & Запись в базу данных, экспорт в холодное хранилище & Асинхронная \\
\bottomrule
\end{tabularx}
\end{table}

\textbf{Эмпирический порог переключения.} Эксперименты показали, что для задач с ожидаемым временем выполнения менее 0.1 мкс синхронное выполнение обеспечивает лучшую совокупную производительность.
При длительности операции свыше 0.5 мкс асинхронная модель становится предпочтительной независимо от типа нагрузки.
В диапазоне 0.1--0.5 мкс выбор зависит от требований к предсказуемости задержки и доступных вычислительных ресурсов.

\subsubsection{Выводы по разделу}

Проведённое сравнительное исследование подтверждает целесообразность разделения стратегий выполнения в системе событий для DPI-приложений:

\begin{enumerate}
    \item Синхронное выполнение обеспечивает минимальную задержку и максимальную предсказуемость для легковесных вычислительных задач, что критично для обработки пакетов в реальном времени.

    \item Асинхронная модель позволяет изолировать блокирующие операции ввода-вывода от основного потока обработки, сохраняя высокую пропускную способность системы при наличии гетерогенной нагрузки.

    \item Гибридная архитектура, поддерживающая обе стратегии с явным разделением на этапе регистрации, позволяет разработчикам принимать осознанные решения о маршрутизации обработчиков на основе характеристик конкретной задачи.

    \item Эмпирический порог в 0.1--0.5 мкс может служить практическим ориентиром для выбора стратегии при проектировании новых модулей анализа.
\end{enumerate}

Таким образом, предложенное архитектурное разделение не является избыточной абстракцией, а представляет собой обоснованный компромисс между производительностью, предсказуемостью и гибкостью, необходимый для построения масштабируемых систем глубокого анализа сетевого трафика.




\newpage

%
%
%
% В рамках проектирования API решаются три ключевые задачи: механизм создания новых типов событий, способ генерации (диспетчеризации) событий и метод регистрации обработчиков.
% Остальные аспекты, включая управление потоками, синхронизацию и распределение памяти, инкапсулируются внутри библиотеки и прозрачны для пользователя.
% Для следования лучшим практикам C++ и предотвращения конфликтов имён вся функциональность инкапсулирована в пространство имён \texttt{event}.
%
% \subsubsection{Создание новых типов событий}
% Система событий оперирует двумя основными сущностями событие (event) и обработчик (handler).
%
% Для удовлетворения требованиям к гибкости необходимо поддержать возможность
%
% Система должна иметь гибкую и простую настройку событий со стороны пользователя, но при этом необходимо вызывать правильные обработчики для каждого события. Для гибкой настройки событий будем использовать расширение через наследование.
%
% Для определения новых типов событий применяется механизм наследования на основе Curiously Recurring Template Pattern (CRTP).
% Данный подход позволяет избежать виртуальных таблиц и RTTI, но при этом гарантировать уникальность идентификатор типа.
% Пользователь определяет структуру события, наследуясь от шаблонного базового класса, который автоматически генерирует уникальный идентификатор типа и предоставляет методы доступа к полям.
%
% \begin{lstlisting}[language=C++, caption={Определение пользовательского типа события}]
% namespace event {
% class IdGenerator {
%  public:
%   static size_t nextId() noexcept {
%     static std::atomic<size_t> id_generator_{0};
%     return id_generator_.fetch_add(1, std::memory_order::relaxed);
%   }
% };
%
% template <typename Derived>
% class EventBase : private IdGenerator {
%  public:
%   static size_t getTypeId() noexcept {
%     static const size_t type_id = nextId();
%     return type_id;
%   }
%
%  protected:
%   ~EventBase() = default;
% };
% } // namespace event
%
% struct PacketEvent : public event::BaseEvent<PacketEvent> {
%     uint32_t protocol;
%     uint16_t size;
%     uint64_t timestamp_ns;
% };
% \end{lstlisting}
%
% Использование CRTP в сочетании с \texttt{IdGenerator} гарантирует, что каждый тип события имеет уникальный линейный идентификатор.
% Линейный рост идентификаторов позволяет эффективно адресоваться к массивам обработчиков за O(1), не требуя использования хеш-таблиц для маршрутизации.
% Поля структуры располагаются в памяти непрерывно, что улучшает локальность кэша и позволяет применять векторизованные операции фильтрации при обработке групп событий.
%
% % TODO: таким образом получаем гибкую систему настройки различных событий
%
% \subsubsection{Генерация события заданного типа}
% Для генерации событий реализуется шаблонная функция \texttt{dispatch}, принимающая произвольное число аргументов через \textit{perfect forwarding}.
% Функция должна обеспечивать типобезопасность на этапе компиляции и эффективное создание объектов событий.
%
% \begin{lstlisting}[language=C++, caption={Механизм диспетчеризации событий}]
% template <typename EventType, typename... Args>
% void dispatch(Args&&... args) {
%     static_assert(std::is_base_of_v<event::BaseEvent<EventType>, EventType>,
%                   "EventType must inherit from BaseEvent");
%
%     // Event dispatch code
% }
% \end{lstlisting}
%
% \subsubsection{Создание обработчиков событий фиксированного типа}
% Регистрация обработчиков осуществляется через шаблонную функцию \texttt{subscribe}, принимающую любой вызываемый объект (лямбда, функция, функтор) с сигнатурой, совместимой с принимаемым типом события.
%
% \begin{lstlisting}[language=C++, caption={Регистрация обработчика событий}]
% template <typename EventType, typename Callable>
% void subscribe(Callable&& callback) {
%     static_assert(std::is_invocable_v<Callable, const EventType&>,
%                   "Callback must accept const EventType&");
%
%     // Handle subscribe code
% }
% \end{lstlisting}
%
% \subsection{Внутренняя реализация и оптимизации}
% \subsubsection{Type Erasure для обработчиков событий}
% Для хранения разнородных обработчиков в едином контейнере применяется техника \textit{type erasure}, позволяющая привести произвольные вызываемые объекты к общему интерфейсу без потери типобезопасности на этапе вызова.
% В отличие от стандартного \texttt{std::function}, в реализации используется реализации \texttt{fu2::unique\_function} из библиотеки \texttt{Naois/function2} \cite{function2_benchmark}.
%
%
% \begin{lstlisting}[language=C++, caption={Handler с type erasure}]
% namespace Core::event {
%
% class Handler {
%  public:
%   template <typename EventType, typename Callable>
%   Handler(std::type_identity<EventType>, Callable&& callback) {
%     static_assert(
%         std::is_invocable_v<Callable, const EventType&>, "Callback must accept const EventType&"
%     );
%     invoke_ = [cb = std::forward<Callable>(callback)](const void* event) {
%       cb(*static_cast<const EventType*>(event));
%     };
%   }
%
%   void handle(const void* event) {
%     assert(invoke_);
%     invoke_(event);
%   }
%
%  private:
%   struct SBOCapacity {
%     static constexpr std::size_t capacity  = 8 * 6;
%     static constexpr std::size_t alignment = alignof(std::max_align_t);
%   };
%
%   template <typename... Signatures>
%   using unique_function = fu2::function_base<true, false, SBOCapacity, true, false, Signatures...>;
%   using InvokeFunction  = unique_function<void(const void*)>;
%
%   InvokeFunction invoke_{};
% };
%
% }  // namespace Core::event
% \end{lstlisting}
%
% Выбор \texttt{fu2::function} вместо \texttt{std::function} обусловлен следующими факторами:
% \begin{itemize}
%     \item \textbf{Расширенный SBO}: стандартный \texttt{std::function} обычно имеет буфер 16--24 байта, что недостаточно для большинства лямбд с захватом нескольких переменных. Настройка буфера в 48 байт покрывает $\sim$95\% типичных обработчиков без динамической аллокации \cite{sbo_performance};
%     \item \textbf{Больше возможностей}: поддержка только перемещаемых функций, поддержка \texttt{const}, \texttt{volatile}, \texttt{noexcept}.
% \end{itemize}
%
% \subsubsection{Реализация хранения и запуска обработчиков}
% Мы будем хранить наши обработчики в такой структуре.
% \begin{lstlisting}[language=C++, caption={}]
% std::vector<std::vector<Handler>> handlers_storage;
% \end{lstlisting}
%
% В ячейке соответствующей идентификатору типа будет лежать массив с обработчиками событий этого типа.
%
% Тогда функция \texttt{subscribe} будет просто сохранять обработчик в этот масив.
%
% \begin{lstlisting}[language=C++, caption={Реализация подписки на событие}]
% template <typename EventType, typename Callable>
% void subscribe(Callable&& callback) {
%   static_assert(std::is_invocable_v<Callable, const EventType&>,
%                         "Callback must accept const EventType&");
%
%   const std::size_t id = EventType::getTypeId();
%   auto& handlers       = detail::handlerStorage().sync;
%
%   if (id >= handlers.size()) {
%     handlers.resize(id + 1);
%   }
%
%   handlers[id].push_back(
%     Handler{std::type_identity<EventType>(), std::forward<Callable>(callback)}
%   );
% }
% \end{lstlisting}
%
% Получаем идентификатор типа события, на которое мы подписываемся, и хранилище обработчиков, расширяем хранилище если нужно и сохраняем обработчик.
%
% \subsubsection{Оптимизация через синхронный запуск}
%
% \newpage
% Архитектура библиотеки строится на принципах разделения ответственности: слой захвата (вне библиотеки), слой нормализации, ядро событийной шины, пулы обработки и слой доставки. Центральным компонентом выступает асинхронный маршрутизатор событий, реализованный на основе направленного ациклического графа (DAG) очередей. Каждый узел графа соответствует типу события или группе обработчиков, а рёбра определяют маршруты доставки с учётом приоритетов и фильтров.
%
% Для обеспечения высокой пропускной способности применяются следующие оптимизации:
% \begin{itemize}
%     \item Lock-free MPMC-очереди на основе атомарных индексов и cache-line выравнивания, исключающие false sharing \cite{wang2021wasm}.
%     \item Пулы памяти с pre-allocated блоками, специфичными для типов событий, что устраняет фрагментацию и ускоряет аллокацию/освобождение.
%     \item Thread-per-core модель с work-stealing для балансировки нагрузки при неравномерном распределении событий.
%     \item Векторизованная фильтрация событий на этапе маршрутизации с использованием SIMD-инструкций для параллельного сравнения полей.
%     \item Backpressure-механизм, адаптирующий скорость генерации к пропускной способности обработчиков через динамическое ограничение диспетчеризации и сброс низкоприоритетных событий при перегрузке \cite{dag2022design}.
% \end{itemize}
%
% Управление конфигурацией и схемами событий вынесено в декларативный реестр, поддерживающий версионирование и атомарные обновления. Это позволяет обновлять логику обработки без остановки конвейера, что критично для систем непрерывного мониторинга. Изоляция обработчиков достигается через разделение контекстов выполнения и ограничение доступа к глобальному состоянию, что предотвращает каскадные отказы при ошибках в пользовательской логике.
%
% \newpage
% TODO:
% Выделение памяти осуществляется из lock-free пула, специфичного для каждого типа события, что устраняет фрагментацию и contention на уровне менеджера памяти.
% Маршрутизация выполняется по типу через предвычисленные индексы, обеспечивая $O(1)$ сложность диспетчеризации.
%
% Для минимизации overhead применяется тип-стирание с опциональной оптимизацией: если обработчик не сохраняет состояние и совместим с указателем на функцию, используется прямой вызов, в противном случае применяется компактный \texttt{std::function}-аналог с small object optimization.
%
% Архитектура гарантирует потокобезопасность регистрации через lock-free атомарные операции и блокировки только на уровне метаданных реестра.
% Обработчики выполняются в контексте worker-потоков с привязкой к ядрам CPU (CPU affinity), что минимизирует переключения контекста и улучшает предсказуемость задержек. Система поддерживает множественных подписчиков на один тип события с настраиваемым порядком вызова (параллельным или последовательным), что позволяет реализовать конвейерную обработку без блокировок.
